\documentclass[11pt]{article}
% Generic article skeleton. The target venue is undecided; replace with the venue's official
% template once it is chosen. This draft has NOT been compiled locally (no LaTeX on the machine).
%
% DRAFT STATE (2026-11-15): only the result-independent parts are written.
% No empirical result exists yet; every place that will carry a result is marked \todo{...}.
% Every claim in this file is listed in paper/claims.md with its source.
\usepackage[margin=1in]{geometry}
\usepackage{amsmath,amssymb,amsthm,mathtools}
\usepackage{graphicx}
\usepackage{booktabs}
\usepackage{xcolor}
\usepackage[numbers]{natbib}
\usepackage{hyperref}

% --- theorem-like environments -------------------------------------------------------------
% IMPORTANT (project rule 3): no result may be called "Theorem" until both external reviewers
% have reviewed the final text with no open issue. Everything is a Proposition / Lemma / Remark.
\newtheorem{proposition}{Proposition}
\newtheorem{lemma}[proposition]{Lemma}
\newtheorem{conjecture}[proposition]{Conjecture}
\theoremstyle{remark}
\newtheorem{remark}{Remark}

% --- macros ----------------------------------------------------------------------------------
\newcommand{\todo}[1]{\textcolor{red}{[TODO: #1]}}
\newcommand{\E}{\mathbb{E}}
\newcommand{\Prob}{\mathbb{P}}
\newcommand{\cF}{\mathcal{F}}
\newcommand{\cG}{\mathcal{G}}
\newcommand{\rentry}{r_T^{\mathrm{entry}}}
\newcommand{\thetasel}{\theta_{\mathrm{sel}}}
\newcommand{\ZH}{Z_H}
\newcommand{\ZO}{Z_O}
\newcommand{\Fsyn}{\cF_t^{\mathrm{syn}}}

\title{What Does a Trade/Abstain Gate Add?\\
A Pre-Specified Evaluation of Selective Trading on A-Share Daily Data\\
under Retrospective Synthetic Fills}
\author{Anonymous}
\date{}

\begin{document}
\maketitle

\begin{abstract}
A trade/abstain gate placed on top of a predictive signal trades less often, and trading less
often by itself changes market exposure and trading costs. An improvement observed after adding a
gate is therefore hard to read as evidence that the gate selects better trades. We study this
question in a long-only, daily, T+1 backtest of Chinese A-shares with transaction costs, in which
all fills are \emph{retrospective synthetic fills} constructed from daily bars. We pre-specify
(i) an estimand for within-batch selection value, the net return per order of a gate in excess of
the exact expectation of a stratified random choice from the same nominated list at matched order
opportunities; (ii) tests of whether two feature groups change that quantity when added to the
gate, namely holder-state variables and observable actions of other identifiable participants;
and (iii) a descriptive, additive split of the profit-and-loss difference between a gated and an
ungated strategy into exposure-schedule, composition and cost terms. As a methodological
complement we define the entry-time loss probability of executed entries, adapt time-uniform
confidence sequences to labels that mature with a delay under arbitrary same-day dependence, and
give a worst-case lower bound on the width of any honest interval. The full inference pipeline is
checked by complete-null simulations before the hold-out period is opened.
\todo{results: outcomes of the pre-specified tests C1--C4, descriptive attribution, monitoring
interval on the synthetic fill stream.}
\end{abstract}

% =============================================================================================
\section{Introduction}
\label{sec:intro}
% =============================================================================================

\paragraph{Motivation.}
A poker player who judges that the hand is weak relative to what the table has revealed folds
and waits. Computer programs that play imperfect-information poker at a superhuman level have
made this style of reasoning prominent \citep{brown2018libratus,brown2019pluribus}. The trading
analogue is immediate: a trader who holds a predictive signal but doubts it on a given day can
decline to act. We use the analogy only as motivation. This paper does not model a game, does not
solve for an equilibrium, and makes no game-theoretic claim.

\paragraph{What existing work already does.}
Declining to act is an established design in quantitative trading. Selective classifiers have
been turned into trading strategies that stay flat when the classifier abstains
\citep{pidan2011selective,chalkidis2021trading,nabar2023conservative}; a strategy-level gate has
been placed on a cross-sectional equity ranker \citep{sanderink2026alpha}; and practitioners use
a secondary model that filters the signals of a primary strategy under the name meta-labeling
\citep{joubert2022metalabeling,meyer2023metalabeling}. The statistical tools for controlling the
error rate on the accepted cases are also available, both for independent or exchangeable data
\citep{geifman2017selective,bates2021rcps,angelopoulos2022crc} and in sequential, anytime-valid
form \citep{howard2021timeuniform,ramdas2023savi,xu2024arcps,khosravi2026csa}. We therefore do
not present the gate architecture, or the idea of abstaining, as a contribution.

\paragraph{Why an improvement from abstention is hard to read.}
A gated strategy differs from its ungated base in three ways at once: how much capital is in the
market on each day, which names are held with that capital, and how much is paid in trading
costs. The difference in additive profit and loss between the two strategies splits exactly into
these three parts (Lemma~\ref{lem:attribution}), and only the second part can reflect better
selection; even that part also contains sector, size and liquidity tilts and the changed mix of
overlapping entry cohorts. A comparison of a gated strategy with its ungated base consequently
does not say whether the gate chose better trades. A second difficulty is statistical: trade
outcomes on the same day are dependent, positions overlap, and outcomes are known only after the
holding period, so the number of trades overstates the amount of evidence
(Proposition~\ref{prop:lower}).

\paragraph{Questions.}
We work in a long-only, daily, T+1 backtest of A-shares with transaction costs, built from free
daily-bar data without opening-auction timestamps. All fills are retrospective synthetic fills
(Section~\ref{sec:ashare}), and every empirical statement in this paper is restricted to that
setting.
\begin{itemize}
\item \textbf{RQ-1 (pre-specified).} Does a trade/abstain gate placed on a multi-factor base
signal add economically non-negligible selection value \emph{within the same nominated batch},
relative to a stratified random choice at matched order opportunities?
\item \textbf{RQ-2 (pre-specified).} Does adding a feature group to the gate change that
within-batch selection value? We examine two groups separately: holder-state variables $\ZH$
(capital-gains overhang, chip-distribution statistics, abnormal turnover) and observable actions
of other identifiable participants $\ZO$ (Dragon--Tiger list disclosures, block trades, executed
insider and large-shareholder share changes; free data only). Neither group is treated as a
measure of anyone's private information.
\item \textbf{RQ-3 (descriptive).} How does the additive profit-and-loss difference between a
gated strategy and the ungated base portfolio split into exposure-schedule, composition and cost
terms, and what are terminal wealth and drawdown at the level of the deployed strategy? No
confirmatory claim is attached to this question.
\end{itemize}

\paragraph{Contributions.}
The contributions below do not depend on how the tests turn out.
\begin{enumerate}
\item \textbf{An evaluation protocol.} A point-in-time, reproducible protocol for evaluating
selective trading on A-share daily data: an explicit definition of synthetic fills and of the
information set they imply, benchmark-price accounting with ledger reconciliation, a random
comparator at matched order opportunities whose expectation is computed exactly, a frozen annual
deployment algorithm, and a complete-null simulation of the whole pipeline that must pass before
the hold-out period is opened (Section~\ref{sec:design}).
\item \textbf{An estimand for selected-trade risk, with a monitoring interval and a lower bound.}
We define the entry-time conditional loss probability of executed entries, state what it is not,
and give a time-uniform interval for it when labels mature with a bounded delay and same-day
outcomes are arbitrarily dependent, together with a worst-case lower bound on the width of any
honest interval (Section~\ref{sec:theory}). The interval is an adaptation of existing tools,
time-uniform confidence sequences \citep{howard2021timeuniform} combined with sub-sampling for
delayed feedback; we do not claim a new type of guarantee. Proofs are in the appendix.
\item \textbf{A complete report of pre-specified tests.} The tests of the two feature groups
$\ZH$ and $\ZO$ in the gate, and of the within-batch selection value of two reference gates, are
specified before any outcome is seen and are reported in full, including negative and
inconclusive outcomes (Section~\ref{sec:design}).
\end{enumerate}
\todo{results: one paragraph summarising the outcomes of C1--C4 and the descriptive findings,
using only the wording fixed in Section~\ref{sec:wording}.}

% =============================================================================================
\section{Related Work}
\label{sec:related}
% =============================================================================================

\paragraph{Abstention as a trading decision.}
The idea that a predictive model may decline to act, and that declining corresponds to holding
no position, has a history in finance. \citet{pidan2011selective} study selective prediction of
financial trends with hidden Markov models. \citet{chalkidis2021trading} turn binary and ternary
selective classifiers into trading strategies that stay flat whenever the classifier abstains;
they construct the selection function with the procedure of \citet{geifman2017selective},
evaluate on 30-minute bars of five CME metal futures with an anchored walk-forward scheme, and
report that abstention raises accuracy on the traded subset in every configuration while most
configurations are unprofitable at 0.3 ticks of slippage. Their guarantee is the one stated for
independent and identically distributed samples; temporal dependence is not discussed.
\citet{nabar2023conservative} train a cascade of models and take no trade when the cascade
refuses to predict, with experiments on the Indian equity market.

More recent work places the abstention decision on top of an existing signal.
\citet{sanderink2026alpha} separates two decisions for a gradient-boosted cross-sectional equity
ranker: a strategy-level gate, built mainly from the trailing realised rank information
coefficient of matured predictions, that decides whether to trade at all on a given date, and a
position-level cap that reduces exposure for the names with the highest predicted rank
displacement. The study uses a universe of up to 100 U.S.\ stocks, reports that the gate predicts
model failure better than volatility-based regime proxies, and provides no formal guarantee.
For the CSI~300 index, \citet{fang2026volatility} report that naive predictive strategies fail
after costs whereas designs with low-volatility gating and threshold calibration improve
defensive performance. Selective prediction has also been applied to cryptocurrency event
contracts \citep{hao2026coverage}, to limit-order-book forecasting \citep{manoharan2026uqlob} and
to retrieval-augmented financial forecasting \citep{torres2026finabstain}. At the strategy
level, \citet{guo2026whennot} studies when a machine-learning strategy for factor and sector
rotation should refuse to trade; according to its abstract, the signal is deployed only when its
validation improvement over a fallback benchmark exceeds a calibrated threshold, and the paper
gives a return and mean--variance utility decomposition and a validation-error deployment bound.
\todo{read the full text of Guo (2026); its decomposition may overlap with
Lemma~\ref{lem:attribution}, in which case the lemma must be presented as following that paper.}

In quantitative-finance practice the same architecture is known as meta-labeling
\citep{joubert2022metalabeling}, and the calibration of the secondary model for position sizing
has been studied by \citet{meyer2023metalabeling}.
\todo{verify the chapter of L\'opez de Prado (2018) that introduces meta-labeling before citing
it as the origin of the term.}
The gates in this paper are, structurally, instances of this line of work. What we add is not a
gate but a way of asking what a gate contributes.

\paragraph{Selective classification and risk control.}
The reject option goes back to \citet{chow1970optimum}, and a learning theory for rejection is
developed by \citet{cortes2016learning}. \citet{geifman2017selective} select a confidence
threshold with a guarantee on the risk of the accepted set, and \citet{geifman2019selectivenet}
integrate the reject option into the network. \citet{ziyin2019deepgamblers} recast selective
classification as a betting problem in which abstaining corresponds to reserving part of the
stake; their theorems characterise the optimal betting solution. \citet{zhang2023abstain},
motivated by finance, posit a data-generating process in which part of the samples is
uninformative and analyse the recovery of the informative region under independent and
identically distributed sampling, without a trading evaluation.

A parallel literature calibrates a threshold after training so that a user-specified risk is
controlled: risk-controlling prediction sets \citep{bates2021rcps}, Learn-then-Test
\citep{angelopoulos2021ltt} and conformal risk control \citep{angelopoulos2022crc}, with
extensions that control risk on the selected cases \citep{xu2025scrc,bai2026score}. For data that
are not exchangeable, adaptive conformal inference addresses distribution shift
\citep{gibbs2021aci}; time-uniform confidence sequences \citep{howard2021timeuniform} and
anytime-valid inference \citep{ramdas2023savi} give guarantees that hold at arbitrary stopping
times; and anytime-valid constructions have been used for risk control on sequentially collected
data \citep{xu2024arcps,hultberg2026anytime} and for selective acting in the deployment of
language models \citep{khosravi2026csa}. Online conformal risk control with delayed feedback is
studied by \citet{wang2026caocrc}, online selective classification with limited feedback by
\citet{gangrade2021online}, and selective classification under distribution shift by
\citet{liang2024selective}.

These results supply the machinery used in Section~\ref{sec:theory}. Our monitoring proposition
should be read as an adaptation of time-uniform confidence sequences to a stream of trades with
delayed labels, overlapping positions and dependent same-day outcomes, not as a new type of
guarantee.
\todo{compare with recent work on anytime-valid inference under delayed outcomes and staggered
entry (not yet read, not yet in the bibliography) before the positioning of
Proposition~\ref{prop:monitor} is final.}

\paragraph{Guarantees at the portfolio level.}
A neighbouring literature attaches guarantees to portfolio-level decisions.
\citet{mohammadpourazari2026safe} let a controller choose among a prediction-independent safe
portfolio and two signal-following overlays, with an evidence gate that moves unsupported
directional mass to the safe portfolio; their statements are a tracking-regret bound under
delayed feedback and a pathwise safety bound, not a statement about the loss probability of
selected trades.
\citet{zarei2026conformal} gate the proposals of a reinforcement-learning allocator with a
conformally calibrated risk budget. Conformal methods have further been used for portfolio
selection \citep{kato2024conformal}, stock selection \citep{kaya2025conformal}, position sizing
\citep{ryan2026conformalkelly} and the calibration of portfolio tail-risk measures
\citep{schmitt2026taming,chen2026cvar}; anytime-valid tests have been used to decide when to
promote a trading rule \citep{benhamou2026promote}. The quantity we monitor is different from
all of these: the entry-time loss probability of executed entries. We make no statement about
drawdown or portfolio tail risk, which depend on position sizing and on the dependence between
overlapping positions.

\paragraph{Information asymmetry and the decision not to trade.}
Why a less informed party might hesitate to trade is a classical topic in economic theory and
market microstructure, including trade under common knowledge \citep{milgrom1982information},
trading against informed counterparties \citep{kyle1985continuous,glosten1985bid}, the pricing of
new issues \citep{rock1986new} and portfolio choice and non-participation under ambiguity
\citep{dow1992uncertainty,easley2009ambiguity}. Whether information risk is priced is itself an
empirical question \citep{easley2002information}. \citet{choi2013informed}, using account-level
holdings data from the Shanghai Stock Exchange, report that stocks whose prior institutional
ownership is more volatile, their predictor of institutional information advantage, earn a
higher four-factor alpha, while the raw return difference is not statistically significant and
the effect is absent among small stocks. This is a caution for any abstention rule: declining
trades in which one suspects a disadvantage may forgo return, which is why our design reports
the return of the orders a gate did not submit (Section~\ref{sec:design}).
\citet{chakraborty2008bluffing} show, in a two-period model with an insider, rational followers
and competitive market makers, that the insider bluffs in every equilibrium; the result concerns
the informed party and does not characterise when the less informed should refrain from trading.
Disclosure of insider trades as a strategic signal is analysed by \citet{huddart2001disclosure}.
Related work on execution and market making treats position building as a game of incomplete
information \citep{chriss2025position} and studies the detection of toxic order flow by brokers
\citep{cartea2026toxic,cartea2025simple}.

\paragraph{The two feature groups.}
The ``chip distribution'' displayed by Chinese retail trading software estimates the
distribution of holders' purchase prices from the history of prices and turnover. Its academic
counterpart is the capital-gains overhang, a reference price weighted by past turnover, which
the literature associates with \citet{grinblatt2005prospect} and with the study of
reference-dependent preferences and the risk--return trade-off by \citet{wang2017reference}.
\todo{both papers are verified at the bibliographic level only; confirm the definition of
capital-gains overhang from the full texts.}
\citet{xu2025replication} replicates the latter study on Chinese data and examines the
interaction between capital-gains overhang and risk proxies, and \citet{pan2025overreaction}
examine capital gains and returns for A-shares. These variables describe the state of current
holders; we call them holder-state variables. The second group consists of publicly disclosed
actions of identifiable participants. Such disclosures have been studied in relation to prices
and returns in the Chinese market: Dragon--Tiger list disclosures \citep{hong2025dragontiger},
block trades \citep{fan2012block} and insider sales \citep{du2024fastsales}; for the U.S.\
market, \citet{lakonishok2001insider} ask whether insider trades are informative. Any
predictive content of these variables is thus not a finding of this paper. We examine only whether adding each group to a gate changes
within-batch selection value, under controls for the price and volume history from which the
variables are built and for the events that trigger disclosure. Machine-learning return
prediction is studied for the Chinese stock market by \citet{leippold2022machine} and in
empirical asset pricing more broadly by \citet{gu2020empirical}; market-specific trading rules have been incorporated into learned
trading systems by \citet{zhang2022t1}, and the predictability of price-limit hits is studied by
\citet{liang2024pricelimit}.
\todo{Chinese-language literature (CNKI, broker research reports) has not been searched; no
statement of the form ``no prior work evaluates these variables in an abstention layer'' is made
until that search and a fact-check are complete.}

% =============================================================================================
\section{Problem Setup}
\label{sec:setup}
% =============================================================================================

This section states the abstract setting on which the statements of Section~\ref{sec:theory}
depend. It contains no market rule, execution model or data source. How A-share trading is mapped
onto it is a separate set of modelling assumptions, summarised in Section~\ref{sec:ashare}.

\subsection{Objects and information times}

Time is indexed by integers $t=1,2,\dots$ (``days''). All random objects live on one probability
space $(\Omega,\mathcal{A},\Prob)$. Each day has two information times. The pre-submission
information $\cG_t$ is what is known when the orders of day $t$ are submitted, before their
fills are known. The post-fill information $\cF_t$ contains $\cG_t$ together with the outcome of
the day-$t$ fills, and $\cG_t\subseteq\cF_t\subseteq\cF_{t+1}$. Orders are chosen with $\cG_t$;
which of them become executed entries is known with $\cF_t$.

On day $t$ there is a finite random candidate set $C_t$. For $i\in C_t$, $O_{ti}\in\{0,1\}$
indicates that an order is submitted and $A_{ti}\in\{0,1\}$ that candidate $i$ is an executed
entry. The number of entries is $n_t=\sum_{i\in C_t}A_{ti}$ and $N_T=\sum_{t\le T}n_t$. Every
executed entry carries a loss label $L_{ti}\in\{0,1\}$; the number of losing entries of day $t$
is $X_t=\sum_{i\in C_t}A_{ti}L_{ti}$, and
\[
\mu_t=\E[X_t\mid\cF_t]
\]
is the entry-time conditional expected number of losses.

\subsection{Assumptions}

\paragraph{A1 (decision information).}
$(\cF_t)_{t\ge1}$ is an increasing family of $\sigma$-fields. For every $t$, the set $C_t$, the
indicators $A_{ti}$ and $n_t$ are $\cF_t$-measurable, and $n_t\le M$ for a known integer
$M\ge1$. Any randomisation used is $\cF_t$-measurable.

\paragraph{A1-G (pre-submission measurability).}
There are $\sigma$-fields $\cG_t\subseteq\cF_t$ and $\cG_t$-measurable order indicators
$O_{ti}$ with $\sum_iO_{ti}\le M$ such that the complete specification of every submitted order
is $\cG_t$-measurable and $A_{ti}\le O_{ti}$. A1-G is used in no proof. It is required in order
to read the results as statements about an implementable strategy: A1 alone would allow a rule
that decides which orders ``participated'' after seeing their fills.

\paragraph{A2 (fixed, frozen label).}
(a) Every executed entry $(t,i)$ has a label $L_{ti}\in\{0,1\}$ defined on all of $\Omega$ by a
convention (execution, valuation, exit and missing-data rules) fixed before any data are
examined. (b) Wherever the convention refers to what data are available, it does so at a
deterministic label-finalisation lag $d_{\mathrm{label}}\ge1$ that is part of the convention,
and the value so obtained is final; in particular the label is not a function of the monitoring
lag $d$ of A3 and is never revised. (c) Nothing is assumed about the joint law of the labels: no
independence or conditional independence within a day, no stationarity, no mixing, and no
restriction on how $\mu_t$ evolves.

\paragraph{A3 (bounded revelation lag).}
There is a known deterministic integer $d\ge1$, fixed before any data are examined, such that
$X_t$ is $\cF_{t+d}$-measurable for all $t\ge1$.

\paragraph{A4 (observable accounting).}
At monitoring time $T$ the monitor knows, for every executed entry $(t,i)$ with $t\le T$,
whether its label has been revealed, and the value of every revealed label, where ``revealed''
means that the final value of the label itself is known. $K_T$ counts the revealed labels equal
to $1$, and $P_T$ is $N_T$ minus the number of revealed labels. No condition is placed on the
revelation mechanism; which labels are revealed early may depend on the labels.

\paragraph{A5 (individual revelation by the lag; optional).}
Every label of every entry of day $t$ is revealed at all monitoring times $T\ge t+d$. Then
$P_T\le dM$. A3 does not imply A5.

\medskip
Each assumption has to be supported in an application by a data audit or a written convention.
A3, in particular, is a substantive assumption about the data pipeline and does not follow from
a holding horizon.

\paragraph{Two lags.}
The label-finalisation lag $d_{\mathrm{label}}$ belongs to the label; the monitoring lag $d$
belongs to the monitoring procedure. Because the labels are frozen by A2(b), a larger
pre-specified monitoring lag leaves $X_t$, $\mu_t$ and the target unchanged. The interval radius
of Proposition~\ref{prop:monitor} is not monotone in $d$, so $d$ must be fixed before the data
are seen; a lag chosen after looking at interval widths is not covered.

\subsection{The target, and what it is not}
\label{sec:target}

The target is
\[
\rentry=\frac{\sum_{t\le T}\mu_t}{N_T},\qquad N_T>0,
\]
the average, over executed entries through $T$, of the loss probability that each day's batch
had given $\cF_t$. It is a random and in general unobservable quantity. It is \emph{not}:
\begin{enumerate}
\item the realised loss fraction $R_T=\sum_{t\le T}X_t/N_T$;
\item the expectation of $R_T$ given the information at time $T$, which would condition old
entries on new information;
\item a statement about the next entry or batch: $\mu_{T+1}$ can be anything given the past;
\item a population parameter: the model contains no stationary ``true loss probability'';
\item independent of the filtration: a different declared $(\cF_t)$ that satisfies A1--A3 gives
a different, also validly monitored, target, so any statement about $\rentry$ must name $\cF_t$;
\item a measure of economic performance: a low loss probability is compatible with rare large
losses.
\end{enumerate}
Monitoring $\rentry$ is inference about the past entry stream. It is not control of future risk.

\subsection{A-share instantiation: modelling assumptions}
\label{sec:ashare}

The mapping of A-share trading onto the abstract objects is a set of modelling assumptions and
conventions. It is not part of any statement in Section~\ref{sec:theory}, and the conclusions of
that section apply to an A-share trade stream only to the extent that these assumptions hold for
it. The statements about exchange rules on which the mapping relies have not yet been checked
against the official rule texts.
\todo{replace this sentence by the outcome of the rule audit once the official texts have been
checked.}
The full specification is in Appendix~\ref{app:exp}; we summarise the parts needed to read the
rest of the paper.

\paragraph{Days, orders and entries.}
A day is a market trading day. Orders are buy orders intended for the opening call auction of
day $t$, at most one per name and at most $M$ per day, generated by code that reads only data
dated up to day $t-1$ and the public pre-open information of day $t$; this is $\cG_t$.

\paragraph{Retrospective synthetic fills.}
Our data are daily bars. A daily bar is produced after the close and does not record whether the
opening auction produced a match or what volume it matched. An entry indicator built from a
daily bar is therefore not measurable with respect to the information available at the auction.
We adopt the following rule and call its output by its name. A buy order for stock $i$ on day
$t$ is recorded as a \emph{synthetic fill}, in full and at the daily open of day $t$, if (S1) the
daily bar of day $t$ exists, has positive volume and is not flagged as a whole-day suspension;
(S2) the daily open is not the upper price limit computed from the reference price; and (S3) the
order size does not exceed a capacity cap computed from $\cG_t$. Otherwise the order is recorded
as not filled and is cancelled. Of the day-$t$ bar the rule reads only the traded flag, the
suspension flag and the open. The declared post-fill information is
\[
\Fsyn=\text{all data dated }\le t-1,\ \text{the pre-open public information of day }t,\
\text{and the day-}t\text{ open, traded flag and suspension flag of all names.}
\]
This is an artificial information set: on a day on which a name had no auction match it contains
facts that became known only later that day. A real order may not have filled where the rule
records a fill, or may have filled at a different price, and the size of this discrepancy cannot
be estimated from daily bars. Accordingly, every monitoring quantity in this paper refers to the
synthetic fill stream with filtration $(\Fsyn)$, and no result is a statement about an
executable strategy.

\paragraph{Label and lags.}
The primary label marks a net loss at the scheduled exit, the opening auction of day $t+H$, with
$H\ge1$ under the T+1 rule. It is a valuation computed by a written decision tree that assigns a
value to every entry, including entries affected by corporate events, suspensions or missing
inputs; unsupported events and missing inputs are recorded as losses and counted. The label is
computed once from the data snapshot of day $t+d_{\mathrm{label}}$ and never revised. Because
the tree uses facts known only at the close of day $t+H$, we fix
$d_{\mathrm{label}}=d=H+1$ before any data are examined. Whether every input is available by
then is an availability assumption to be audited, not an established fact; with a single
downloaded file it can hold only as a function of that file.

\paragraph{Accounting.}
Holdings, returns, cost bases and marks of all strategies use the same benchmark price at each
instant, the daily open. Explicit charges and modelled execution shortfall enter only through
the cost term, each exactly once (Lemma~\ref{lem:attribution}).

% =============================================================================================
\section{Monitoring, Limits of Evidence, and Attribution}
\label{sec:theory}
% =============================================================================================

\paragraph{Review status and attribution.}
The statements of this section have been examined in three rounds by two external model
reviewers, which reported no error in the proofs or constants of
Propositions~\ref{prop:monitor}--\ref{prop:lower} or in the identity of
Lemma~\ref{lem:attribution}; the revisions made after the last round have not been re-reviewed.
Pending that confirmation we label the results propositions.
\todo{update after the confirmatory review; rename only if both reviewers report no open issue.}
Proposition~\ref{prop:monitor} combines two existing ingredients, a stitched time-uniform
boundary of the kind developed by \citet{howard2021timeuniform} and the reduction of delayed
feedback to interleaved sub-sequences known from online learning with delays
\todo{cite the delayed-feedback sub-sampling reduction (Weinberger--Ordentlich 2002; Joulani et
al.\ 2013) once the entries are verified and added to the bibliography.}
We do not claim it as a new type of result. Proofs are deferred to Appendix~\ref{app:proofs}.

\subsection{A time-uniform interval under delayed labels}

Fix $\delta\in(0,1)$ and a deterministic integer $d\ge1$ for which A3 holds; both are fixed
before any data are examined. Let $c_0=\pi^2/6$ and, for integers $k\ge0$ and real $q\ge0$,
\[
L_k=\log\frac{2dc_0(k+1)^2}{\delta},\qquad
g_k(q)=\tfrac12\sqrt{2^kL_k}\,\Bigl(1+\frac{q}{2^{k+1}}\Bigr).
\]
Days are partitioned into $d$ lanes by the residue of the day index modulo $d$. For lane
$j\in\{0,\dots,d-1\}$ let
\[
Q_{j,T}=\sum_{\substack{t\le T\\ t\equiv j\ (\mathrm{mod}\ d)}}n_t^2 ,
\]
a non-negative integer. For integer arguments define $b(0)=b^*(0)=0$ and, for $q\ge1$,
\[
b(q)=\sqrt{2^{k(q)}L_{k(q)}},\quad k(q)=\lfloor\log_2q\rfloor,\qquad
b^*(q)=\min_{k\ge0}g_k(q),
\]
and set $B_T=\sum_{j=0}^{d-1}b(Q_{j,T})$ and $B_T^*=\sum_{j=0}^{d-1}b^*(Q_{j,T})$.

\begin{proposition}[Delayed monitoring under arbitrary within-day dependence]
\label{prop:monitor}
Assume A1--A3. There is an event $E_\delta$ with $\Prob(E_\delta)\ge1-\delta$ on which,
simultaneously for all lanes $j$ and all $T\ge1$,
\begin{equation}
\Bigl|\sum_{\substack{t\le T\\ t\equiv j\ (\mathrm{mod}\ d)}}(\mu_t-X_t)\Bigr|
\le b^*(Q_{j,T})\le b(Q_{j,T}).
\label{eq:lane}
\end{equation}
Consequently, under A4 in addition (so that $K_T$ and $P_T$ are computable), on $E_\delta$ and
simultaneously for all $T$ with $N_T>0$:
\begin{enumerate}
\item[(a)] $\displaystyle\max\Bigl\{0,\frac{K_T-B_T}{N_T}\Bigr\}\le\rentry\le
\min\Bigl\{1,\frac{K_T+P_T+B_T}{N_T}\Bigr\}$;
\item[(b)] $\displaystyle\max\Bigl\{0,\frac{K_T-B_T^*}{N_T}\Bigr\}\le\rentry\le
\min\Bigl\{1,\frac{K_T+P_T+B_T^*}{N_T}\Bigr\}$, and $B_T^*\le B_T$;
\item[(c)] $E_\delta=E_\delta^{\mathrm{up}}\cap E_\delta^{\mathrm{low}}$ with
$\Prob(E_\delta^{\mathrm{up}})\ge1-\delta/2$ and $\Prob(E_\delta^{\mathrm{low}})\ge1-\delta/2$;
the upper bounds in (a) and (b) hold on $E_\delta^{\mathrm{up}}$ and the lower bounds on
$E_\delta^{\mathrm{low}}$.
\end{enumerate}
When $N_T=0$ the target is undefined and no claim is made. The interval is valid for any
revelation pattern consistent with A3--A4, including revelation times that depend on the labels.
\end{proposition}
% STATUS: Theorem 1(a)(b)(c) of theory/theorems.md. `GPT 审过`（半独立）、`Gemini 审过`（Flash）.
% 三轮均未被判错；第 3 轮后的修改未复审。审计认为基本是已知工具的推论，放附录；
% 在日线合成成交下只对合成成交流成立。 Author's own label in theorems.md: `自查通过`.
% (c): no separate numerical check of the one-sided levels.
% NOT to be called a theorem until the confirmatory review reports no open issue.

The cap $M$ of A1 does not enter the radii; it makes $X_t$ bounded, so that $\mu_t$ is an
ordinary conditional expectation. The proposition is a statement for each pair $(\delta,d)$
fixed in advance and does not cover a lag selected after looking at the data. It is a
Hoeffding-type bound with a predictable range clock, $\sum n_t^2$, and dyadic stitching, applied
on $d$ interleaved sub-sequences. It is not sharper than known confidence sequences and has no
empirical-variance adaptation.

The idea is the following. Within one lane, the outcome of every day is known before the next
day of that lane is traded, so the difference between expected and realised loss counts behaves
like a fair game inside the lane and cannot drift far; summing the $d$ lanes gives a band around
the observed loss count. Entries whose label is not yet revealed are counted both ways, as no
loss in the lower endpoint and as a loss in the upper endpoint.

\begin{proposition}[Rate of the sampling-error term; total width]
\label{prop:rate}
Let $\ell_T=\log(2dc_0/\delta)+2\log\bigl(1+\log_2(MN_T)\bigr)$.
\begin{enumerate}
\item[(a)] Under A1--A3, for every $T$ with $N_T>0$, deterministically,
\[
B_T^*\le B_T\le\sqrt{d\,\ell_T\sum_{t\le T}n_t^2}\le\sqrt{dMN_T\,\ell_T}.
\]
Under A4 in addition, the total width of each interval in
Proposition~\ref{prop:monitor}(a),(b) is at most
$\min\bigl\{1,\ P_T/N_T+2\sqrt{dM\ell_T/N_T}\bigr\}$.
\item[(b)] If A5 also holds, then $P_T\le dM$ and the total width is at most
\[
\min\Bigl\{1,\ \frac{dM}{N_T}+2\sqrt{\frac{dM\ell_T}{N_T}}\Bigr\}\le3\sqrt{\frac{dM\ell_T}{N_T}}.
\]
\end{enumerate}
\end{proposition}
% STATUS: Corollary 1.1 of theory/theorems.md. 同 Theorem 1：`GPT 审过`（半独立）、`Gemini 审过`（Flash）；
% 第 3 轮后的修改未复审。非贡献。 Author's label: `自查通过`. Part (b): no dedicated numerical check.

Under A1--A4 alone nothing is said about $P_T/N_T$, because A3 concerns the daily loss count and
not the disclosure of individual labels. The total width is therefore not claimed to be of order
$\sqrt{dM/N_T}$ without A5; it tends to zero only if both $dM\ell_T/N_T\to0$ and
$P_T/N_T\to0$.

\begin{proposition}[Realised loss fraction]
\label{prop:realised}
For all $T$ with $N_T>0$, under A4 alone and without any probabilistic assumption,
\[
\frac{K_T}{N_T}\le R_T\le\frac{K_T+P_T}{N_T}.
\]
Under A1--A3, on the event $E_\delta$ of Proposition~\ref{prop:monitor},
$|\rentry-R_T|\le B_T^*/N_T$ for all such $T$.
\end{proposition}
% STATUS: Corollary 1.2 of theory/theorems.md. 同 Theorem 1：`GPT 审过`（半独立）、`Gemini 审过`（Flash）；
% 第 3 轮后的修改未复审。非贡献。 Author's label: `自查通过`.

All rates above use the common denominator $N_T$, the number of executed entries. The
``matured-only'' rate $K_T/(N_T-P_T)$ has a different denominator and a composition selected by
maturity; it is covered by no statement here and can be biased. For example, with two entries
whose labels are $0$ and $1$, of which only the losing one is revealed, $R_T=1/2$ and the
bracket is $[1/2,1]$, while the matured-only rate equals $1$.

% ---------------------------------------------------------------------------------------------
% DRAFT, NOT REVIEWED -- not part of the paper's statements.
% Corollary 1.3 of theory/theorems.md (interval for a target label with unbounded resolution
% time, e.g. executable liquidation, via certified bracket labels; assumptions A1, A2(a), CERT,
% A3 for the two bracket labels with lag d_D, A4-D).
% STATUS: `草稿`. 第 3 轮 Gemini 判正确、GPT 给出反例；已改用认证指示并新增假设 CERT，未复审。
% If it survives review it may be stated as a Proposition; until then it may at most appear as a
% Conjecture, and no empirical statement may rely on it unconditionally.
% ---------------------------------------------------------------------------------------------
An extension to labels whose resolution time has no deterministic bound, such as a label based
on executable liquidation, is in draft and has not passed review; it is not stated here and
Proposition~\ref{prop:monitor} is not applied to such labels.

\subsection{How much evidence a trade stream can carry}

The sampling-error term of Proposition~\ref{prop:rate} shrinks like $\sqrt{dM/N}$ rather than
$\sqrt{1/N}$. The next statement shows that, in the worst case over environments satisfying the
assumptions, this dependence cannot be improved.

\paragraph{Block environment.}
Let $d,M,B\ge1$ be integers, $T=Bd$ and $N=MBd$. Days $1,\dots,T$ are grouped into $B$ blocks of
$d$ consecutive days, every day has exactly $M$ entries, and there are no entries after $T$. Let
$Z_1,\dots,Z_B$ be independent Bernoulli($p$) variables; every label of every day in block $b$
equals $Z_b$. For $t\le T$, $\cF_t$ is generated by the $Z$'s of the blocks before that of day
$t$ and by an independent uniform variable $U$ available for randomisation. Call the law
$\Prob_p$. Under $\Prob_p$, A1--A5 hold with lag $d$, $\mu_t=Mp$ for $t\le T$, and $\rentry=p$.
A pair of measurable functions $\ell\le u$ of $(Z_1,\dots,Z_B,U)$, that is, of everything that
is ever observable in this environment, is called a $(1-\delta)$-honest interval at time $T$
over a family of laws if $\Prob(\ell\le\rentry\le u)\ge1-\delta$ for every law in the family.

\begin{proposition}[Worst-case effective sample size]
\label{prop:lower}
Let $\delta\in(0,1/2)$ and $h=\frac{\sqrt3}{8}(1-2\delta)\sqrt{Md/N}$. If $(\ell,u)$ is
$(1-\delta)$-honest at time $T$ over the two-point family
$\{\Prob_{1/2-h},\Prob_{1/2+h}\}$, then for each $p\in\{1/2-h,\,1/2+h\}$,
\[
\Prob_p\Bigl(u-\ell\ge\tfrac{\sqrt3}{4}(1-2\delta)\sqrt{Md/N}\Bigr)\ge\frac{1-2\delta}{2},
\qquad\text{hence}\qquad
\E_p[u-\ell]\ge\tfrac{\sqrt3}{8}(1-2\delta)^2\sqrt{Md/N}.
\]
Moreover, in this environment
$\operatorname{Var}\bigl(\sum_{t\le T}(X_t-\mu_t)\bigr)=p(1-p)\,MdN$, against $p(1-p)N$ for
independent labels.
\end{proposition}
% STATUS: Proposition 2 of theory/theorems.md. `GPT 审过`（半独立）、`Gemini 审过`（Flash）.
% 三轮均未被判错；**新颖性未验证**（须先读 Weinberger–Ordentlich 2002）. Author's label: `自查通过`.
% Round-3 revisions not re-reviewed.

Any procedure that is honest over all environments satisfying A1--A4 with parameters $(d,M)$ is
honest over this two-point family, and the interval of Proposition~\ref{prop:monitor} at a fixed
$T$ is one such procedure. Since A5 holds in the block environment,
Proposition~\ref{prop:rate}(b) bounds the width of that interval by
$3\sqrt{dM\ell_T/N}$. Over the class of environments satisfying A1--A5, the dependence of the
width on $dM/N$ therefore cannot be improved in the worst case beyond the factor
$\sqrt{\ell_T}$ and constants: in the block environment $N$ entries carry the information of
$N/(Md)$ independent ones. The argument is a standard two-point comparison, and the construction
with outcomes constant on blocks of length $d$ resembles lower-bound constructions for online
learning with delayed feedback.
\todo{verify novelty against Weinberger--Ordentlich 2002}
We make no claim about whether this bound was previously known.

The proposition has clear limits. It says nothing about the logarithmic factors or about
$\delta\to0$. It is a worst case over the class and does not say that a given market resembles
the block environment; on a restricted class, for example with bounded cross-sectional
dependence, narrower honest intervals may exist, but that restriction would be an additional
assumption that data on one path cannot certify. It is stated only for $N$ a multiple of $Md$,
its constants are not optimised, and it concerns $\rentry$, not the realised fraction $R_T$,
which is bracketed deterministically by Proposition~\ref{prop:realised}.

\subsection{An additive attribution identity}

Let periods be $t=1,\dots,T$. For $a\in\{G,B\}$ (gated, baseline) let $h_t^a\in\mathbb{R}^m$ be
the vector of risky-asset holdings in currency units carried through period $t$,
$x_t\in\mathbb{R}^m$ the vector of asset returns over period $t$ in excess of the common cash
return, $C^a\in\mathbb{R}$ the total trading cost, and define the additive performance quantity
\[
\Pi^a=\sum_t(h_t^a)^\top x_t-C^a .
\]
Let $e_t^a\ge0$ and $w_t^a\in\mathbb{R}^m$ be any numbers and vectors with
$h_t^a=e_t^aw_t^a$; for a long-only strategy, $e_t^a$ is the invested capital and $w_t^a$ the
portfolio weights.

\begin{lemma}[Additive attribution identity]
\label{lem:attribution}
With the above definitions,
\begin{equation}
\Pi^G-\Pi^B
=\underbrace{\sum_t(e_t^G-e_t^B)(w_t^B)^\top x_t}_{\text{exposure-schedule term}}
+\underbrace{\sum_te_t^G(w_t^G-w_t^B)^\top x_t}_{\text{composition term}}
-\underbrace{(C^G-C^B)}_{\text{cost term}} .
\label{eq:attribution}
\end{equation}
Moreover, for the exposure-schedule reference $h_t^E=e_t^Gw_t^B$ with any assigned cost $C^E$
and $\Pi^E=\sum_t(h_t^E)^\top x_t-C^E$, the gross difference
$\sum_t(h_t^G-h_t^E)^\top x_t$ equals the composition term,
$\Pi^G-\Pi^E$ equals the composition term minus $(C^G-C^E)$, and $\Pi^E-\Pi^B$ equals the
exposure-schedule term minus $(C^E-C^B)$.
\end{lemma}
% STATUS: Lemma 3 of theory/theorems.md. `GPT 审过`（半独立）、`Gemini 审过`（Flash）.
% 恒等式两方确认；记账解释经三轮修改，末轮未复审；不主张新颖性；与 Guo 2026 的分解是否重合未知。
% Author's label: `自查通过` (algebra; numerically exact to rounding).

The identity is pathwise algebra and needs no assumption. It is an accounting definition of the
same kind as textbook performance attribution into allocation, selection and interaction,
applied over time to exposure; we claim no novelty for it.
\todo{cite a verified source for classical performance attribution (believed to be Brinson, Hood
and Beebower 1986; not yet opened or in the bibliography).}
Several qualifications govern its use.
\begin{enumerate}
\item \emph{Composition is not stock picking.} The composition term contains every difference
between $w^G$ and $w^B$: sector, size and liquidity tilts and the mix of overlapping entry
cohorts as well as selection. This is why the pre-specified tests of Section~\ref{sec:design}
use a different, within-batch quantity.
\item \emph{Conventions.} When $e_t^B=0$ the weights $w_t^B$ are not determined by $h_t^B$, and
the split between the first two terms depends on a convention that must be fixed in advance. The
decomposition is not unique: with the reference $e_t^Bw_t^G$ the two components change by the
interaction term $\sum_t(e_t^G-e_t^B)(w_t^G-w_t^B)^\top x_t$, which \eqref{eq:attribution} books
under composition.
\item \emph{The reference need not be feasible.} $h^E$ may violate trading constraints. It is an
attribution reference and is not an admissible comparator for the value of a deployed strategy.
\item \emph{Costs.} The identity requires one return vector $x_t$ for all strategies. We
therefore record all strategies at a common benchmark price at each instant and book explicit
charges and execution shortfall in $C^a$, so that every economic cost is counted exactly once.
\item \emph{Wealth.} $\Pi^a$ equals the change in wealth of strategy $a$ under sufficient
accounting conditions (self-financing, total returns measured consistently with holdings, all
costs counted once, zero cash return); otherwise it is reported as additive excess profit and
loss as defined.
\item \emph{Scope.} Sharpe ratio, maximum drawdown, expected shortfall and compounded return are
not linear in the profit-and-loss path and do not inherit \eqref{eq:attribution}. The identity
attributes realised numbers; it does not identify a causal effect of timing or selection.
\end{enumerate}

\subsection{Monitoring is not control}

\begin{remark}[No guarantee of successful trading in every environment]
\label{rem:nocontrol}
Call an environment in which every candidate label equals $1$ almost surely the all-loss
environment, and let $\alpha<1$ and $T\ge1$. In the all-loss environment, for every policy,
$R_T=\rentry=1$ on $\{N_T>0\}$. Hence no policy can guarantee, in every environment satisfying
A1--A4, a positive probability of trading with risk at most $\alpha$, in either sense
$\rentry\le\alpha$ or $R_T\le\alpha$; and any policy that guarantees risk at most $\alpha$
almost surely on $\{N_T>0\}$ in every such environment must abstain almost surely in the
all-loss environment.
\end{remark}
% STATUS: Remark 4 of theory/theorems.md. `GPT 审过`（半独立）、`Gemini 审过`（Flash）.
% 陈述在第 1 轮后改弱。第 3 轮后的修改未复审。 Author's label: `自查通过`. Elementary; no novelty claimed.

The remark does not say that the only universally safe policy is never to trade, nor that
entry-time information cannot be informative. Its content is that a guarantee of the form ``risk
at most $\alpha$ together with a positive probability of trading'' requires restrictions on the
environment or a feasibility assumption. For this reason the statements of this section concern
monitoring of the past entry stream and never a promise about future losses.

% =============================================================================================
\section{Experimental Design}
\label{sec:design}
% =============================================================================================

This section describes the design only. Nothing in it is a result. All parameter values quoted
are those of the pre-specified plan.
\todo{replace planned values by the frozen values and give the hash of the frozen conventions
file.}
All fills are retrospective synthetic fills under the rule of Section~\ref{sec:ashare}, and
every table, figure and sentence that reports an outcome carries that qualifier.

\subsection{Data, universe and timeline}

The backtest uses free daily-bar data for A-shares.
\todo{data source, access date, licence, and the outcome of the data audit (point-in-time
constituents, delisted names, corporate actions, tradability flags).}
The planned universe is the point-in-time union of the constituents of the CSI~300 and CSI~500
indices. The planned development segment ends on 2021-12-31 and is used for tuning and for all
design-stage calculations; the planned hold-out period, 2022-01 to 2026-06, is frozen and opened
once.
\todo{confirm universe and periods after the data audit.}
Models, calibration maps and thresholds are refitted once a year by a frozen deployment
algorithm that uses only earlier data, with a purge gap of $H+1$ trading days between training
labels and deployment.

\subsection{Base signal, order opportunities and the standard order}

A multi-factor base predictor built from price and volume features ranks the candidate set each
day, and its top $K$ names form the \emph{nominated list}. In all pre-specified tests the
nominated list comes from the base predictor without either feature group, so that all gates
choose from the same lists. A gate decides, from $\cG_t$, which nominated names to submit.

For every nominated name $i$ on day $t$ let $\rho_{ti}$ be the net profit and loss, per unit of
slot cash, of one \emph{standard order}: a pre-determined quantity from $\cG_t$, a synthetic
fill under the rule of Section~\ref{sec:ashare}, a scheduled exit at the open of day $t+H$ with
re-submission after exits recorded as not filled, benchmark-price accounting with the full cost
model, and $\rho_{ti}=0$ if the order is recorded as not filled. $\rho_{ti}$ depends on the name
and the day only, not on the gate.

\subsection{Estimands}
\label{sec:estimands}

\paragraph{Within-batch selection value.}
The nominated list of each day is divided into strata by liquidity and volatility terciles
computed from data through day $t-1$. For gate $g$ let $\mathrm{Sub}_t^g$ be the names it
submits, $S_t^g=|\mathrm{Sub}_t^g|$, $S_{t,s}^g$ its count in stratum $s$, and
$\bar\rho_{t,s}$ the mean of $\rho$ over all nominated names of stratum $s$. Define
\[
u_t^g=\sum_{i\in\mathrm{Sub}_t^g}\rho_{ti}-\sum_sS_{t,s}^g\,\bar\rho_{t,s},
\qquad
\thetasel^g=\frac{\sum_tu_t^g}{\sum_tS_t^g}.
\]
The second sum in $u_t^g$ is the exact expectation of the first under a uniformly random choice
of $S_{t,s}^g$ names from each stratum, so no Monte-Carlo comparator is needed. $\thetasel^g$ is
the order-weighted mean net return of the gate's orders in excess of a stratified random choice
from the same nominated list, at matched order opportunities: the same number of orders, the
same cash per order and the same per-stratum counts on each day. It does not match the invested
exposure of deployed strategies. By construction it gives no credit for avoiding whole days:
$u_t^g=0$ on days on which the gate submits all or none of the list. Because cells from which no
different choice of the same size was possible contribute zero, we report beside every
$\thetasel^g$ the share of orders that come from cells with a real choice and the range that
the statistic could attain given the gate's own per-cell counts.

\paragraph{Feature-group increments.}
For the holder-state group, $\psi_H=\thetasel^{\mathrm{LOSS}(+\ZH)}-\thetasel^{\mathrm{LOSS}(-Z)}$,
and for the group of observable actions of other identifiable participants,
$\psi_O=\thetasel^{\mathrm{LOSS}(+\ZO)}-\thetasel^{\mathrm{LOSS}(-Z)}$, where the gates are those
of Table~\ref{tab:methods} and choose from the same nominated lists.

\paragraph{Descriptive quantities.}
The value of deployed strategies is described by the annualised arithmetic return, terminal
wealth, drawdown, trade-level loss magnitude, turnover, and the return of orders that a gate did
not submit. Comparisons of deployed strategies use feasible strategies only, with the same
initial capital, execution rules, cost model and dates. The difference in additive return
between a gated strategy and the ungated base is split by Lemma~\ref{lem:attribution} into
exposure-schedule, composition and cost terms. The interval of Proposition~\ref{prop:monitor} is
reported over time for the synthetic fill stream with filtration $(\Fsyn)$; it describes how
much evidence the stream can carry and is never used to compare methods.

\subsection{Gates, baselines and feature groups}

\begin{table}[t]
\centering
\caption{Methods of the first round. ``Within-batch'' means that the method can choose a strict
subset of the nominated list, so that $\thetasel$ is informative for it. All methods use the
same information cut-offs, deployment algorithm, tuning budget and selection criterion.}
\label{tab:methods}
\small
\begin{tabular}{@{}llp{0.52\linewidth}l@{}}
\toprule
Name & Type & Definition & Within-batch \\
\midrule
G0 & reference & submit all $K$ nominated orders & no \\
SCORE & gate & submit the top $\lceil cK\rceil$ names by base score & yes \\
CONF & gate & submit if the calibrated probability of a positive net return exceeds a global
threshold & yes \\
LOSS($-Z$) & gate & gradient-boosted classifier of the loss label on base and control features;
submit if the predicted loss probability is below a threshold & yes \\
LOSS($+\ZH$), LOSS($+\ZO$) & gate & the same class, grid and budget with one feature group
added & yes \\
LOSS-LR($-Z$) & gate & regularised logistic version of LOSS($-Z$) & yes \\
NET-K, NET-U & gate & regression of the standard order's net return, on the nominated list
(NET-K) or on the full candidate set (NET-U) & yes \\
VOL, LIQ & filter & drop names with high volatility or low traded amount & yes \\
TREND & filter & submit the full list or nothing according to a market-trend rule & no \\
PERF & gate & submit the full list or nothing according to the matured rank information
coefficient of the base score, after \citet{sanderink2026alpha} & no \\
VT & scaling & volatility-targeted scaling of all orders, never above one & no \\
E-const & scaling & constant scaling of G0 to a gate's mean invested fraction, in a feasible
version (previous year) and an ex-post reference version & no \\
E-sched & reference & exposure-schedule reference $h^E=e^Gw^B$ of
Lemma~\ref{lem:attribution}; need not be feasible & --- \\
RND & comparator & stratified random selection at matched order opportunities, simulated only
for path-dependent descriptive metrics & --- \\
\bottomrule
\end{tabular}
\end{table}

Table~\ref{tab:methods} lists the methods. CONF and LOSS($-Z$) are in-house implementations in
one pipeline; statements about them are not statements about existing abstention methods as a
class. Published selective-prediction methods
\citep{geifman2017selective,geifman2019selectivenet,ziyin2019deepgamblers,zhang2023abstain,khosravi2026csa,xu2024arcps}
are left to an optional second round. Equal numbers of tuning trials are not equal optimisation
effort, and we report the search space, the number of distinct configurations and the tuning
cost of every method.

\paragraph{Feature groups.}
$\ZH$, the holder-state variables, comprises capital-gains overhang
(Section~\ref{sec:related}), chip-distribution statistics computed by a turnover-decay method,
and abnormal turnover. $\ZO$, the
observable actions of other identifiable participants, comprises stock-day aggregates of public
disclosure events: Dragon--Tiger list disclosures, block trades and executed insider and
large-shareholder share changes. No participant identity enters any feature. Each group is
always accompanied by control features: the price and turnover history over the windows from
which $\ZH$ is constructed, and the price, range and turnover conditions that trigger the
disclosures in $\ZO$. A class of $\ZO$ enters the analysis only if it passes a data audit,
including an audit of announcement dates, before any modelling.

\paragraph{Factorial structure.}
The pre-specified contrasts add one group to the gate while the base predictor, and hence the
nominated lists, are fixed. Cells in which a group enters the base predictor, both groups enter
the gate, or the group's information is routed to the gate through a score are exploratory. A
capacity-matched control gate receives the columns of a group after a within-day permutation
across names, so that the number of features and their marginal distributions are preserved
while stock-specific content is removed.

\subsection{Pre-specified tests C1--C4}
\label{sec:tests}

All four tests are evaluated once, on the hold-out period, under the primary parameter values
($H=5$ trading days, target coverage $c=0.5$, series averaged over five seeds).
\todo{confirm $K$, $c$, $H$ and the cost parameters as frozen.}
\begin{itemize}
\item \textbf{C1.} $\thetasel^{\mathrm{CONF}}$ lies within the equivalence bound
$\pm\Delta_{\mathrm{ord}}$. Equivalence test with null
$|\thetasel^{\mathrm{CONF}}|\ge\Delta_{\mathrm{ord}}$.
\item \textbf{C2.} The same for $\thetasel^{\mathrm{LOSS}(-Z)}$.
\item \textbf{C3.} $\psi_H=0$ against a two-sided alternative.
\item \textbf{C4.} $\psi_O=0$ against a two-sided alternative. C4 is included only if at least
one class of $\ZO$ passes the data audit.
\end{itemize}

\paragraph{Equivalence bound.}
The bound is anchored to capital: $\Delta_{\mathrm{ord}}=\Delta_{\mathrm{cap}}(H+1)/(Ac)$, where
$\Delta_{\mathrm{cap}}$ is an additive annual return difference per unit of capital fixed in
advance and $A=242$ is the annualisation constant. The adopted value
$\Delta_{\mathrm{cap}}=1\%$ per year corresponds to about 5 basis points per order at $H=5$ and
$c=0.5$. The bound is fixed before any data are loaded and is never changed afterwards.

\paragraph{Status of C1 and C2.}
Whether C1 and C2 are run as confirmatory equivalence tests or reported as pre-declared
estimation targets is decided by a mechanical rule from a design-stage variance calculation that
uses no output of any gate: they are confirmatory only if the bound is at least the smallest
bound with 80\% power under a reference model of a selector without skill, evaluated for the
most persistent reference selector considered. As estimation targets they keep the same
estimand, comparator and bootstrap, the point estimate and intervals are reported, and no
verdict is issued. The status is recorded before any gate is fitted on the development
evaluation years and is not revisited.
\todo{state the recorded status of C1 and C2 and the resulting family size.}

\paragraph{Multiple comparisons.}
The confirmatory tests form one family, adjusted by the Holm step-down procedure at family-wise
level 5\%.
\todo{cite verified references for the Holm procedure, the two one-sided tests procedure and the
stationary bootstrap; none is in the bibliography yet.}
The family has between one and four members, depending on the status of C1 and C2 and on
whether C4 is included; its size is fixed before any gate outcome is seen. Everything else,
including all comparisons of deployed strategies, all attribution terms, all other gates and all
monitoring quantities, is exploratory or descriptive and is reported with unadjusted intervals.

\paragraph{Inference.}
The unit of resampling is the trading date, so that within-day dependence is preserved. We use a
paired stationary block bootstrap over dates applied jointly to all strategies, with ratio
statistics recomputed as ratios of resampled sums and a mean block length proportional to
$H+1$. Seed variation is reported separately and is never merged with a bootstrap interval. The
block bootstrap assumes approximate stationarity and weak dependence beyond the block length,
which regime shifts and rule changes violate; whether the procedure attains its nominal error
rates on data of this kind is not assumed but checked, as described next.

\subsection{Complete-null simulation of the pipeline}

Before the hold-out period is opened, the entire pipeline, comprising tuning, calibration,
thresholding, the matched-opportunity statistic, the bootstrap and the Holm step, is re-run
inside each replication of simulated null worlds built from the development years, with the
nominated lists held fixed.
\begin{itemize}
\item \emph{Outcome permutation.} Outcomes are permuted among names within cells defined by day,
nomination status and stratum. This sets $\thetasel$, $\psi_H$ and $\psi_O$ to zero in
expectation but breaks the dependence between a name's outcomes on neighbouring days.
\item \emph{Feature swap.} For each calendar year an independent random bijection of names is
drawn, and the whole feature row that a gate sees, including the base score as a gate input, is
replaced by that of the partner name, while the real base score still determines the nominated
lists. Outcomes and their dependence are untouched.
\item \emph{Group permutation.} Only the columns of one feature group are permuted across names
within each day, so that the reference gate remains informative. In this world the increment is
not zero by construction, because a gate with additional columns is a different fitted model;
the check is centred on the mean over replications.
\item \emph{Boundary versions.} The bound $\Delta_{\mathrm{ord}}$ is added to the return of
every submitted order, which places $\thetasel$ on the boundary of the equivalence nulls.
\end{itemize}
The pass criterion, fixed in advance, limits the number of replications with a false rejection
by an exact binomial count, and a failure triggers a fixed ladder of at most three
progressively more conservative inference procedures, each evaluated on fresh replications and
all reported. The arguments that the first two worlds have zero selection value by construction
are our own and have not been independently reviewed; the simulation includes a direct check of
them. Passing the criterion does not establish that error rates are at most nominal; it fails to
find rates clearly above nominal. The replications share one realised price history.

\subsection{Sensitivity analyses and wording}
\label{sec:wording}

Two alternative fill determinations, one stricter and one looser than the primary rule, are
applied to the frozen models and orders. They bound the dependence of the outcomes on
conditions (S2) and (S3); all three rules are retrospective, so they do not measure the gap
between synthetic and real fills. Further pre-specified analyses, all exploratory, vary the
holding period, the coverage target, the cost parameters and the tuning budget.

The wording of outcomes is fixed in advance. ``Economically negligible (within
$\pm\Delta_{\mathrm{ord}}$)'' is used only when an equivalence null is rejected after
adjustment; ``a difference was detected'' only when a difference null is rejected after
adjustment; in every other case the outcome is ``no detectable difference; inconclusive''. An
estimate outside the bound is reported descriptively and is not a verdict. A detected increment
of a feature group is described as a change in within-batch selection in this pipeline given the
controls, not as evidence about private information or about a game. Every verdict sentence
names the gate as implemented here, the universe and period, and the synthetic execution model.

\subsection{Results}

\todo{results: data audit and fill rates by reason.}

\todo{results: design-stage variance calculation; recorded status of C1 and C2; family size.}

\todo{results: complete-null simulation, all attempts with counts.}

\todo{results: C1--C4 on the hold-out period, with choice share, attainable range, fill
sensitivity, and the gate-based power report.}

\todo{results: descriptive outcomes of deployed strategies and the three attribution terms;
return of orders not submitted.}

\todo{results: monitoring interval of Proposition~\ref{prop:monitor} on the synthetic fill
stream; dependence diagnostics.}

\todo{results: exploratory factorial cells and robustness analyses; negative and anomalous
results.}

% =============================================================================================
\section{Limitations}
\label{sec:limitations}
% =============================================================================================

\paragraph{Synthetic fills.}
All fills are retrospective synthetic fills constructed from daily bars. The declared
information set is artificial, a real order may not have filled where a fill is recorded, and
the difference between synthetic and real opening-auction fills cannot be estimated without
auction data, which we do not have. No result of this paper is a statement about an executable
strategy, and the fill sensitivity analysis does not close this gap.

\paragraph{One market, one history, free data.}
The study covers one market, one universe and one period, with one base signal, long-only, at
daily frequency. All resampling and all null simulations are conditional on a single realised
price history. The data are free vendor data; a downloaded file is a snapshot as of the download
date, so point-in-time properties hold at most as a function of that file, and the completeness
of the vendor's corporate-event calendar is an audit question.
\todo{report which audit items passed and which did not.}
The group $\ZO$ is restricted to disclosures available in free data; scheduled share unlocks and
pre-announced selling plans are outside the scope.

\paragraph{Unverified operational assumptions.}
The mapping of A-share trading to the abstract setting, comprising the execution model, the
valuation rule, the lags and the treatment of distributions, consists of modelling assumptions.
The exchange rules on which it relies have not yet been checked against the official rule texts.
The propositions apply to the trade stream only to the extent that these assumptions hold.

\paragraph{Status of the theoretical statements.}
The statements of Section~\ref{sec:theory} have been reviewed by two external models, one of
which took part in the design discussions and one of which is not that provider's strongest
model; the revisions after the last round have not been re-reviewed, and we therefore call the
results propositions. The extension to labels with unbounded resolution time is a draft on which
the reviewers disagreed, and it is not stated. Proposition~\ref{prop:monitor} is an adaptation
of known tools. The novelty of Proposition~\ref{prop:lower} has not been verified, and recent
work on anytime-valid inference with delayed outcomes has not yet been compared with our
setting.
\todo{update after the confirmatory review and the remaining literature checks.}

\paragraph{What the monitored quantity does not give.}
The interval concerns the past entry stream of the synthetic fills and a target that depends on
the declared filtration. It is not a guarantee about future trades, about drawdown or about
portfolio tail risk, and a low loss probability is compatible with rare large losses. The
interval is wide: by Proposition~\ref{prop:lower} its sampling-error term cannot shrink faster
than $\sqrt{dM/N}$ in the worst case, and a design-stage evaluation of the formulas, made before
any data were seen, puts its half-width at the end of the planned hold-out period at roughly
$0.17$--$0.22$. It is therefore unsuitable for comparing methods and is not used for that
purpose.
\todo{replace the design-stage estimate by the realised width.}

\paragraph{Limits of the estimands and of the inference.}
$\thetasel$ gives no credit for avoiding whole days, so the value of market timing can appear
only in the descriptive quantities. The composition term of Lemma~\ref{lem:attribution} is not
stock picking, the split between exposure and composition depends on stated conventions, and the
attribution is not causal. The tests of C1 and C2 concern two in-house gates and do not
generalise to abstention methods as a class. The equivalence bound may exceed what the design
can resolve, in which case C1 and C2 are reported as estimates without a verdict; inconclusive
outcomes are possible for every test and are reported as such. The block bootstrap relies on
approximate stationarity, and the null simulation checks the inference procedure on one history
only.

\paragraph{Literature coverage.}
Chinese-language sources and broker research have not been searched, several cited works were
read at the level of the abstract only, and the full text of one closely related preprint
\citep{guo2026whennot}, whose decomposition may overlap with Lemma~\ref{lem:attribution}, has
not been read. We therefore make no claim of priority for any component of this paper.
\todo{update after the supplementary search and the citation check.}

\paragraph{Unfavourable results.}
\todo{results: list every setting in which a gate or a feature group performed worse than its
comparator, and every anomalous or failed run.}

% =============================================================================================
\section{Conclusion}
\label{sec:conclusion}
% =============================================================================================

\todo{results: conclusion, written after the hold-out analysis, using the wording rules of
Section~\ref{sec:wording}.}

\bibliographystyle{plainnat}
\bibliography{references}

\appendix

\section{Proofs}
\label{app:proofs}
% Written by the researcher in a separate file; not edited here.
\IfFileExists{appendix_proofs.tex}{\input{appendix_proofs}}{\todo{appendix\_proofs.tex (complete
proofs of Propositions~\ref{prop:monitor}--\ref{prop:lower}, Lemma~\ref{lem:attribution} and
Remark~\ref{rem:nocontrol}) is not yet available.}}

\section{A-Share Instantiation and Experimental Details}
\label{app:exp}
% Written by the researcher in a separate file; not edited here.
\IfFileExists{appendix_experiments.tex}{\input{appendix_experiments}}{\todo{appendix\_experiments.tex
(A-share modelling assumptions, valuation decision tree, cost model, tuning details, licences of
data and software) is not yet available.}}

\end{document}
